\renewcommand{\arraystretch}{1.3}
\begin{tabularx}{\textwidth}{@{}p{4cm}X@{}}
\toprule
\textbf{Kategorie} & \textbf{Zweck und Logik} \\
\midrule
SIEM & Zentrale Sammlung, Korrelation und Analyse von Logdaten zur Erkennung sicherheitsrelevanter Ereignisse. Unterstützt bei der Identifikation von Angriffsmustern und Anomalien. \\
\addlinespace
EDR & Überwachung und Analyse von Aktivitäten auf Endgeräten zur frühzeitigen Erkennung und Reaktion auf Bedrohungen. Nutzt Verhaltensanalysen und Telemetriedaten. \\
\addlinespace
NAD & Netzwerkbasierte Anomalieerkennung durch Analyse von Traffic-Flüssen und Kommunikationsmustern. Erkennt ungewöhnliche oder verdächtige Netzwerkaktivitäten. \\
\addlinespace
Firewall & Kontrolle und Überwachung des ein- und ausgehenden Datenverkehrs. Erkennt und blockiert bekannte Angriffsmuster oder verdächtige Aktivitäten anhand von Signaturen oder Heuristiken. \\

\addlinespace
Forensik & Analyse und Aufarbeitung von Sicherheitsvorfällen. Dient der Ursachenforschung, Beweissicherung und strukturierten Reaktion auf Angriffe. \\
\addlinespace
Täuschung & Einsatz von Täuschungstechniken wie Honeypots oder Ködern zur frühzeitigen Erkennung von Angreifern. Liefert Warnsignale bei unautorisierten Zugriffen. \\
\bottomrule
\end{tabularx}
\caption{Übersicht sicherheitsrelevanter Technologien und deren Funktionslogik}
\label{tab:sicherheitstechnologien}
